\chapter{Developer Guide: Contributing to DAGonStar}
\label{app:developer}

This appendix is addressed to contributors and agents modifying DAGonStar
itself, rather than to workflow authors; it is reproduced here so that
readers of this textbook who wish to extend the engine --- for instance, by
adding a new task type alongside \texttt{LLM}, \texttt{WEB}, and
\texttt{NATIVE} --- have the project's own development discipline available
as a reference.

\section{Why a Textbook Includes a Developer Guide}

It is worth pausing to justify, briefly, why this appendix exists at all in
a book primarily addressed to workflow \emph{authors} rather than engine
\emph{contributors}. Two reasons motivate its inclusion. First, a reader who
encounters one of this book's many documented limitations
(\Cref{sec:known-constraints} being the most concentrated instance) and
wishes to address it directly --- rather than merely working around it ---
needs exactly the procedural knowledge this appendix provides. Second, and
more broadly, reading how a project's own contributors are asked to extend
it is one of the most reliable ways to understand a system's actual design
boundaries, often more reliably than reading the system's own marketing or
high-level architecture description: a project's testing and extension
discipline reveals, concretely, which parts of the system its own
maintainers consider stable and trustworthy, and which parts they
explicitly flag as requiring caution. The seven-step new-task-type recipe
below, for instance, is the same recipe that produced every task type this
book treats as a stable, well-documented mechanism in Part~II.

\section{Ground Rule: Implementation First}

Documentation, examples, and tests must describe the \emph{current}
implementation in the repository. If intended behavior differs from current
behavior, the project's policy is to either implement and test it, or label
it explicitly as future work --- the same fidelity principle this textbook
has followed throughout.

This ground rule is worth connecting explicitly to the architectural-honesty
discussion of \Cref{ch:architecture} (\Cref{sec:known-constraints}): the
documented constraints reproduced there exist, precisely, \emph{because}
this ground rule was followed when that section of the documentation was
written --- a project that instead documented intended, aspirational
behavior would have no structural reason to mention incomplete shell
quoting or partial test coverage at all. The rule is therefore not merely a
contributor instruction; it is the generative principle behind much of what
makes the documentation this book is built from trustworthy as source
material in the first place.

\section{Development Setup}

\begin{lstlisting}[style=bash]
git clone https://github.com/DagOnStar/dagonstar.git
cd dagonstar
python3 -m venv .venv
. .venv/bin/activate
python -m pip install --upgrade pip
python -m pip install -e .
\end{lstlisting}

Install optional integration dependencies only when needed:

\begin{lstlisting}[style=bash]
python -m pip install -e ".[docker]"
python -m pip install -e ".[cloud]"
python -m pip install -e ".[globus]"
python -m pip install -e ".[api]"
python -m pip install -e ".[all]"
\end{lstlisting}

Note the use of \texttt{-e} (editable install) here, in contrast to the
plain \texttt{pip install -r requirements.txt} path given for workflow
\emph{authors} in \Cref{ch:getting-started}: an editable install links the
installed package directly to the working source tree, so that changes made
to \texttt{dagon/} source files take effect immediately on the next
interpreter invocation, without requiring a reinstall --- the correct mode
for someone actively modifying the engine, as opposed to someone merely
authoring workflows against an unmodified installation.

\section{Test Commands}

\begin{lstlisting}[style=bash, caption={Run unit tests.}]
python -m unittest discover -s tests -v
\end{lstlisting}

\begin{lstlisting}[style=bash, caption={Compile source and tests.}]
python -m py_compile dagon/*.py dagon/api/*.py dagon/communication/*.py dagon/ftp_publisher/*.py dagon/peer2peer/*.py tests/*.py
\end{lstlisting}

\begin{lstlisting}[style=bash, caption={Lightweight quality checks.}]
python -m ruff check dagon/shell.py tests
python -m mypy dagon/shell.py
\end{lstlisting}

(Install \texttt{mypy==1.13.0} and \texttt{ruff==0.6.9} first if not already
available.) When packaging metadata changes:

\begin{lstlisting}[style=bash]
python -m pip install -e .
\end{lstlisting}

\section{Test Strategy}

The current unit tests deliberately avoid external services. They verify
stable core behavior: configuration parsing; workflow defaults; dependency
discovery; cycle detection; serialization; runtime-secret validation.

When adding features:

\begin{enumerate}[leftmargin=1.6em]
  \item add a unit test for local deterministic behavior;
  \item mock external services where useful;
  \item keep Docker/Slurm/SSH/Globus/cloud tests separate from generic CI
    unless the environment is explicitly provisioned;
  \item document manual integration verification commands in the change
    notes.
\end{enumerate}

This four-point strategy is, point for point, the same mock-service
discipline already generalized in \Cref{ch:examples-catalog}'s discussion of
mock services as a verification strategy: just as a workflow \emph{author}
verifies an \texttt{LLM}- or \texttt{WEB}-task mechanism locally against a
mock before relying on a real provider, an engine \emph{contributor}
verifies new engine behavior locally, against mocked external services,
before that behavior is exercised against a real Docker daemon, SSH host,
or Slurm scheduler. The same underlying principle --- separate verifying the
mechanism from verifying the integration --- operates at both the workflow-
authoring layer and the engine-development layer.

\begin{practicebox}
The test suite is intentionally a starting point, not a complete quality
gate. It includes checkpoint reuse, staging command generation, command
failure propagation, packaging extras, and mocked Docker, cloud, Globus, and
API-extra boundaries. CI applies Ruff to the tests and shell helper, and
mypy to the shell helper; both scopes are expanded gradually as modules gain
annotations. Priority areas for further expansion are command rewriting for
\texttt{workflow://} references, local batch integration workflows, and
mocked SSH and Slurm boundaries.
\end{practicebox}

\section{Adding a New Task Type}

\begin{enumerate}[leftmargin=1.6em]
  \item Add a value to \texttt{TaskType}.
  \item Add a mapping in \texttt{tasks\_types}.
  \item Implement a concrete task class.
  \item Ensure the class initializes the base \texttt{Task} correctly.
  \item Override \texttt{on\_execute()} or other environment-specific
    hooks.
  \item Add unit tests for construction and dependency behavior.
  \item Update \texttt{reference\_guide.md}, \texttt{architecture.md}, and
    relevant tutorials (in this book's terms: Appendix~A,
    \Cref{ch:architecture}, and the corresponding lesson chapter).
\end{enumerate}

This is precisely the process by which \texttt{LLM}, \texttt{WEB}, and
\texttt{NATIVE} were added to the original \texttt{BATCH}/\texttt{SLURM}/
\texttt{DOCKER}/\texttt{CLOUD} backend family --- the resulting uniformity
across all seven concrete task types, exploited throughout Part~II of this
book, is a direct consequence of every task type following the same
seven-step recipe.

\begin{example}
Consider, as a thought experiment grounded entirely in this recipe, what
adding a hypothetical \texttt{TaskType.QUEUE} (a task that publishes a
message to a message-queue system rather than waiting for a synchronous
HTTP response, in contrast to \texttt{TaskType.WEB}'s synchronous request
model) would require: a new \texttt{TaskType.QUEUE} enum value (step~1); a
new mapping to, say, a \texttt{dagon.queue.QueueTask} class (step~2); the
class itself, implementing message construction and publication (step~3);
correct base-class initialization preserving \texttt{prevs}/\texttt{nexts}
dependency bookkeeping (step~4); an \texttt{on\_execute()} override
implementing the actual publish call (step~5); unit tests verifying that a
\texttt{workflow://} reference inside the message payload is correctly
parsed and staged, mirroring the existing \texttt{WebTask} tests (step~6);
and updates to Appendix~A, \Cref{ch:architecture}'s task-type table, and a
new lesson chapter analogous to \Cref{ch:web-tasks} (step~7). No part of
this hypothetical extension would require modifying \texttt{Workflow},
\texttt{Task.pre\_run()}, or the \texttt{workflow://} parser itself ---
exactly the architectural payoff the factory-pattern discussion of
\Cref{sec:factory-pattern-depth} predicted.
\end{example}

\section{Modifying Dependency Resolution}

Dependency resolution is implemented primarily in \texttt{Task.pre\_run()}
and \texttt{Task.pre\_process\_command()} (\Cref{ch:schema}). Be careful
with: parsing \texttt{workflow://} references; same-workflow versus
transversal references; reference counts used by garbage collection;
recursive calls to \texttt{workflow.make\_dependencies()} for parallel
modes; service-backed task lookup when the DAGon service is enabled.

Any change here should include tests for: same-workflow references; missing
references; cycles; command rewriting; reference counts.

\section{Modifying Staging}

Staging logic lives in \texttt{Stager.stage\_in()} and
\texttt{Stager.generate\_command()} (\Cref{sec:data-staging-architecture}).
Risks: shell quoting; remote/local path distinction; external command
availability; differences between files, directories, and glob patterns;
service-specific credentials.

\begin{practicebox}
Prefer adding helpers and unit tests before broad rewrites. New staging or
command-generation code should avoid raw string concatenation for paths or
user-controlled values; prefer explicit quoting helpers and tests covering
spaces, special characters, missing files, directories, and glob patterns.
\end{practicebox}

This guidance connects directly to the whitespace boundary rule already
encountered as a workflow-author-facing constraint in
\Cref{sec:whitespace-rule}: the parser limitation that forces a workflow
author to avoid spaces in referenced paths is, from the contributor's side
of the same boundary, exactly the kind of shell-quoting risk this section
warns against introducing further instances of. A contributor strengthening
staging-path quoting is, in effect, working to relax the constraint that
\Cref{sec:whitespace-rule} currently asks workflow authors to observe
manually.

\section{Security Expectations}

\begin{itemize}[leftmargin=1.4em]
  \item Never add real credentials to source files, docs, sample config,
    tests, or CI.
  \item Use placeholders in examples.
  \item Read runtime secrets from environment variables or local ignored
    files.
  \item Avoid logging secrets.
  \item Quote shell arguments defensively in new code.
\end{itemize}

These expectations are the implementation-level counterpart of the
configuration-level rules given in \Cref{ch:configuration} and exercised
throughout the LLM- and Web-task chapters (\Cref{ch:llm-tasks,ch:web-tasks}).
Note the symmetry: \texttt{api\_key\_env}, \texttt{token\_env}, and the
other environment-variable-indirection mechanisms documented from a
workflow author's perspective in \Cref{ch:configuration} exist specifically
\emph{because} contributors are held, by this section's expectations, to
never hard-code credential-handling logic that would make such indirection
unnecessary or bypassable.

\section{Documentation Expectations}

When changing implementation: update the relevant documentation file in
\texttt{docs/}; update \texttt{README.md} if the change affects onboarding;
update \texttt{AGENTS.md} if contributor policy changes; add or update
tests.

\begin{practicebox}
Scientific documentation should be precise enough for a researcher to
reproduce the computational method, and honest about operational
assumptions --- the same standard this textbook has attempted to meet
throughout, including its explicit reproduction of every documented
limitation and open technical-debt item.
\end{practicebox}

\section{Current Technical Debt}

\begin{itemize}[leftmargin=1.4em]
  \item Automated tests cover core local behavior, checkpointing, staging,
    command failure propagation, packaging extras, and selected mocked
    optional-boundary behavior; broader integration coverage is still
    needed.
  \item Public APIs have broad annotations across workflow, task,
    configuration, staging, task-subclass, and API-client code; type
    coverage continues to be tightened when touching public methods.
  \item Shell command construction uses centralized POSIX quoting for core
    task launchers, Slurm, Docker, remote lifecycle commands, staging
    paths, and SCP remote-endpoint prefixes; it should still be extended
    when touching legacy/site-specific integrations.
  \item Optional integrations now have package extras, while
    \texttt{requirements.txt} remains a full development/demo environment;
    lazy imports and explicit extras for Docker, cloud, Globus, API, and
    other service-specific requirements continue to be preserved.
  \item Some legacy examples require site-specific data or services
    (\Cref{ch:examples-catalog}).
\end{itemize}
